\documentclass{article}

\title{Code Review}
\author{}
\date{}

\begin{document}

\maketitle

\section{Review of GD Algorithm for DDT}

The algorithm initializes the S-box with $\texttt{S[0] = 0}$ and also sets $\texttt{S[1]} = \min R_1$.
Computing the minimum of $R_1$ is done in $\mathcal{O}(n)$ time, where $n = 2^k$ is the size of the S-box.

There is a precomputation involved for the algorithm, which calculates $R_i$ for all $i$ from 1 to $n$.
This is done because the algorithm needs to iterate over these sets during the $\texttt{ComputePossibleValues}$ subroutine later.
Calculating $R_i$ for all $i$ takes $\mathcal{O}(n^2)$ time.
It is done by simply iterating through the DDT and finding the non-zero entries.

Now we compute the S-box values for $i$ from 2 to $n$ using a backtracking approach in the $\texttt{RecursiveSearch}$ subroutine.

In each call to $\texttt{RecursiveSearch}$, we calculate the possible values for $\texttt{S[i]}$ using the $\texttt{ComputePossibleValues}$ subroutine.
The subroutine, which iterates through the sets $R_j$ for $j < i$ to find possible values for $\texttt{S[i]}$.
This takes $i \cdot |R_i|$ time, for each $i$, leading to a total of $\mathcal{O}(n^2)$ time for each call.

After calling the $\texttt{ComputePossibleValues}$ subroutine, the search explores all possible values for $\texttt{S[i]}$ and backtracks if all entries are filled and the DDT is not satisfied.
Asymptomatically, the backtracking can lead to an exponential ($\mathcal{O}(n!)$) time complexity in the worst case, especially if the DDT has many non-zero entries, which leads to a large search space.
The exact performance of the algorithm cannot be determined without specific details about the DDT and the structure of the sets $R_i$.

Asymptomatically, the algorithm performs equally well for finding a single solution or all solutions, as it explores the entire search space in both cases.
However, in practice, finding all solutions may take significantly more time than finding just one solution, especially if there are many valid S-boxes that satisfy the DDT.

\section{Review of Enhanced GD Algorithm for DDT}

The Dynamic Table-Optimized Search (DTOS) algorithm works similarly to the GD algorithm but improves the number of guesses.
Again, asymptomatic analysis of the algorithm is not really helpful.
The algorithm performs same as the GD algorithm in terms of time complexity, but it may significantly reduce the search space in practice, leading to faster execution times.

Every time, we fix a value for $\texttt{S[i]}$, or we backtrack and remove the value of $\texttt{S[i]}$, we update the DDT.
The DDT update is done in $\mathcal{O}(n)$ time, as we need to iterate through the DDT entries corresponding to $\texttt{S[i]}$.
The DDT update can significantly reduce the search space, as it eliminates values that would violate the DDT constraints.

\end{document}